\begin{abstract}
LLM 智能体（agent）要在多轮（turn）中行动（调用工具、检索、写入记忆），但训练时通常只有一个结果奖励（outcome reward），于是强化学习必须判断哪些轮该得到功劳。GRPO 及其 2025--2026 年的各种变体给出的答案是“所有轮，一视同仁”：它们改进的是如何优化一个对整条回复统一的优势（advantage），而不是由谁来获得它。本文综述 LLM 智能体的信用分配（credit assignment），重点是 2025--2026 年的工作。我们把方法归为六大类（过程奖励、蒙特卡洛与树式分支、锚状态（anchor state）分组、轮级奖励与 critic、事后（hindsight）与反事实信用、复兴的价值模型），并按两点刻画每一类：它对环境有什么要求，额外算力花在哪里。我们细读了 Memory-R2：其 LoGo-GRPO 从共享的记忆状态出发做重 rollout（rerollout），以修正不公平的组内基线，但它相对 GRPO 的增益缺少等算力（matched compute）对照。一项关于数学工具调用中轮级信用的预注册研究给出三点教训：非守恒的信用实际上充当了奖励，会被策略钻空子（过度调用工具）；守恒约束消除了这一病态，但在不超过四轮的设置下没有可检测的增益；DAPO 式过滤在终点上的增益是用 2.53$\times$ 的 rollout 换来的，在等算力下随之消失。这一领域方法众多，却几乎没有办法衡量信用分配得对不对。我们的首要建议是：构建一个 horizon（时域）可控、真实信用可计算的基准，并与等算力训练挂钩。
\end{abstract}
